\section{Hai chỗ phép đo đã chạy được}
\label{sec:ch18-hai-cho-da-do-duoc}

\index{khoảng trống nghiên cứu!hai chỗ đã đo được|(}%
Như vậy khoảng trống không phải chuyện thiếu số liệu. Mục này nói nó là chuyện
gì, và cách nói là đặt cạnh nhau hai chỗ mà phép đối chiếu đã chạy được, rồi hỏi
hai chỗ ấy có gì chung.

\index{khoảng trống nghiên cứu!chỗ thứ nhất, chuỗi suy luận}%
Chỗ thứ nhất là bài neo, và
chương~\ref{ch:chi-so-khong-do-do-trung-thuc} đã đọc nó hết. Điều mục này cần
không phải kết quả mà là cách bài lấy được vế phải: nó chọn những nhiệm vụ mà đáp
án đúng chỉ ra được nếu mô hình có đi qua một bước nào đó, nên với những bước ấy
thì biết trước bước nào phải xảy ra, không cần nhìn vào trong mô hình. Có vế phải
rồi thì chấm được từng chỉ số, và tám chỉ số cho kết quả quanh mức ngẫu nhiên.

\index{khoảng trống nghiên cứu!chỗ thứ hai, nút thắt khái niệm}%
Chỗ thứ hai là bài mà mục~\ref{sec:ch15-dung-cu-duoc-kiem} đọc, và nó là chỗ mà
chương~\ref{ch:nut-that-khai-niem} gọi là dụng cụ đo đầu tiên đã được kiểm định
mà sách gặp. Ở đó vế phải rẻ hơn nữa: các khái niệm đã được gán nhãn sẵn cùng bộ
dữ liệu, nên không ai phải dựng gì. Có nhãn rồi thì đối chiếu được hai điểm số
\tn{rò rỉ}{leakage} với một tiêu chí mà chính tác giả không chọn ra, trên tám bộ
dữ liệu.

\index{khoảng trống nghiên cứu!cái hai chỗ ấy có chung}%
Đặt hai chỗ ấy cạnh nhau thì cái chung hiện ra, và nó không phải cái mà một người
đọc đoán trước. Cả hai không thông minh hơn phần còn lại của lĩnh vực. Cả hai
đơn giản là đã ở sẵn trong tình huống mà vế phải có mà không phải trả giá cho nó:
một bên vì thiết kế nhiệm vụ ép nó lộ ra, một bên vì người làm dữ liệu đã gán nhãn
từ trước cho một mục đích khác. Không bên nào dựng ground truth cho việc kiểm
định chỉ số; cả hai nhặt được nó.

Hình~\ref{fig:ch18-hinh-dang-khoang-trong} xếp ba phía cạnh nhau, hai phía đã đo
được và một phía chưa.

\begin{figure}[htbp]
  \centering
  \begin{tikzpicture}[
  every node/.style={font=\footnotesize, align=center},
  side/.style={draw=black!85, line width=1pt, fill=black!8,
    minimum width=34mm, minimum height=8mm},
  gap/.style={draw=black!55, dotted, thick, fill=white,
    minimum width=34mm, minimum height=8mm},
  cell/.style={draw=black!65, rounded corners=2pt, fill=black!4,
    minimum width=34mm, minimum height=13mm, text width=32mm,
    font=\scriptsize},
  cellgap/.style={draw=black!55, dotted, thick, fill=white,
    rounded corners=2pt, minimum width=34mm, minimum height=13mm,
    text width=32mm, font=\scriptsize},
  flow/.style={-{Stealth[length=2mm]}, thick, black!70},
  lbl/.style={font=\scriptsize, align=center, black!55}
]
  \node[side] (a) at (0,3.1)   {chuỗi suy luận};
  \node[side] (b) at (4.1,3.1) {nút thắt khái niệm};
  \node[side] (c) at (8.2,3.1) {attribution đặc trưng};

  \node[cell] (a1) at (0,1.4)   {ground truth từ \textbf{thiết kế nhiệm vụ}: đáp án đúng buộc phải đi qua một bước};
  \node[cell] (b1) at (4.1,1.4) {ground truth từ \textbf{nhãn có sẵn}: khái niệm đã được gán nhãn cùng dữ liệu};
  \node[cellgap] (c1) at (8.2,1.4) {không ai biết mô hình đã dùng gì; mọi chỉ số dùng một đại lượng thay thế};

  \node[cell] (a2) at (0,-0.6)   {phép đối chiếu \textbf{đã chạy}: tám chỉ số, kết quả quanh mức ngẫu nhiên};
  \node[cell] (b2) at (4.1,-0.6) {phép đối chiếu \textbf{đã chạy}: hai điểm số, tương quan trên tám bộ dữ liệu};
  \node[cellgap] (c2) at (8.2,-0.6) {phép đối chiếu \textbf{chưa ai chạy}};

  \draw[flow] (a.south) -- (a1.north);
  \draw[flow] (b.south) -- (b1.north);
  \draw[flow] (c.south) -- (c1.north);
  \draw[flow] (a1.south) -- (a2.north);
  \draw[flow] (b1.south) -- (b2.north);
  \draw[flow] (c1.south) -- (c2.north);

  \node[lbl] at (0,-2.05)   {chương~\ref{ch:chi-so-khong-do-do-trung-thuc}};
  \node[lbl] at (4.1,-2.05) {chương~\ref{ch:nut-that-khai-niem}};
  \node[lbl] at (8.2,-2.05) {Phần~II và Phần~IV};

  \node[lbl, text width=30mm] at (11.4,1.4)  {vế phải của định nghĩa~\ref{def:ch01-do-trung-thuc} lấy ở đâu};
  \node[lbl, text width=30mm] at (11.4,-0.6) {đã đem chỉ số ra đối chiếu với nó chưa};
\end{tikzpicture}

  \caption{Hàng giữa là chỗ quyết định. Ở hai cột trái, vế phải của
  định nghĩa~\ref{def:ch01-do-trung-thuc} có sẵn vì một lý do không liên quan gì
  tới việc kiểm định chỉ số, nên hàng dưới chạy được. Ở cột phải, nơi mọi phương
  pháp của Phần~II sống, hàng giữa để trống, và hàng dưới để trống theo. Nét chấm
  mang nghĩa thiếu, như từ chương~\ref{ch:lime-nao-dang-tin} tới giờ.}
  \label{fig:ch18-hinh-dang-khoang-trong}
\end{figure}

\index{khoảng trống nghiên cứu!phía không nhặt được thì chưa ai trả giá}%
Cột thứ ba là phía mà cả Phần~II lẫn Phần~IV bàn tới, và nó là phía duy nhất
trong ba phía mà không ai nhặt được vế phải. Ở đó, muốn biết mô hình đã dùng gì
thì phải nhìn vào trong nó, và
mục~\ref{sec:ch16-ba-loi-thoat} đã ghi rằng nội dung tính toán bên trong không
quan sát trực tiếp được. Cho nên phía ấy phải \emph{dựng} lấy vế phải, tức phải
trả giá, và chương~\ref{ch:ground-truth-dung-san} đã đếm được bốn cách dựng cùng
cái giá của từng cách.

Bây giờ đọc lại cả ba cột như một câu. Ở hai chỗ mà giá bằng không thì phép đối
chiếu đã được chạy. Ở chỗ mà giá khác không thì nó chưa được chạy. Vậy điều phân
biệt hai bên không phải là ai nghĩ ra việc phải kiểm định dụng cụ đo của mình:
bài siêu đánh giá của mục trước nghĩ ra từ năm 2023 và nói rõ vì sao không làm
được. Điều phân biệt là ở hai bên kia có người đã trả sẵn cái giá ấy vì một lý do
khác.

\index{khoảng trống nghiên cứu!cách dựng ground truth không gắn với đối tượng}%
Còn một câu nữa phải nói ở đây, vì thiếu nó thì cột thứ ba trông như một chỗ
không đi vào được. Cách dựng ground truth của bài neo không gắn với
\tn{chuỗi suy luận}{chain-of-thought}.
Mục~\ref{sec:ch09-doi-tuong-khac-nhau} đã ghi điều đó từ chín chương trước,
và ghi nó vào danh sách những thứ đi qua được ranh giới: nguyên tắc chọn những
bài toán mà đầu ra để lộ phép tính đã sinh ra nó là một nguyên tắc về thiết kế
nhiệm vụ, không phải về loại lời giải thích. Attribution đặc trưng dùng được
nguyên tắc ấy. Cái nó phải trả là công dựng bộ nhiệm vụ, và cái nó nhận lại là
một phạm vi hẹp mà mục~\ref{sec:ch16-gia-duong-dung-san} đã đo trên chính bài
neo.

\index{khoảng trống nghiên cứu!hai chỗ đã đo được|)}%
Nên khoảng trống, đọc tới đây, đã hẹp lại thành một câu cụ thể hơn câu mà
mục~\ref{sec:ch09-doi-tuong-khac-nhau} để lại. Không phải không ai đo được. Là
chưa ai chịu trả giá để đo, ở đúng cái phía mà giá khác không, trong khi ở hai
phía giá bằng không thì việc ấy đã được làm hai lần và cho hai câu trả lời khác
nhau. Mục sau xếp lại những điều kiện mà sách đã gom được cho một dụng cụ như
vậy, nếu có ai chịu trả.
